\documentclass[11pt]{article}
\usepackage[letterpaper,margin=0.92in,top=0.85in,bottom=0.82in,headheight=15pt]{geometry}
\usepackage{amsmath,amssymb,amsthm,booktabs,array,graphicx,xcolor,lmodern,fancyhdr}
\usepackage[colorlinks=true,linkcolor=teal!55!black,urlcolor=teal!55!black,citecolor=teal!55!black]{hyperref}
\usepackage{enumitem}
\definecolor{ink}{HTML}{173947}
\newcommand{\heading}[1]{\par\noindent{\large\bfseries\color{ink}#1}\par\smallskip}
\newcommand{\labeltext}[1]{\textbf{\color{ink}#1}}
\newcommand{\QQ}{\mathbb Q}
\newcommand{\CC}{\mathbb C}
\newcommand{\RR}{\mathbb R}
\newcommand{\Cl}{\operatorname{Cl}}
\newcommand{\e}{\mathrm e}
\newcommand{\dd}{\,\mathrm d}
\newcommand{\Li}{\mathcal L}
\newcommand{\ord}{\operatorname{ord}}
\newcommand{\up}{https://github.com/lsernau/grh-lean/blob/72fead102f663d5c7a407964a13dc0cf9d977207/grh}
\newcommand{\gh}[2]{\href{\up/#1}{#2}}
\pagestyle{fancy}\fancyhf{}\lhead{\small\sffamily GRH AND ITS ARITHMETIC CONSEQUENCES}
\rhead{\small\sffamily Research note}\cfoot{\small\thepage}
\renewcommand{\headrulewidth}{0.3pt}
\setlength{\parindent}{0pt}\setlength{\parskip}{5pt plus 0.5pt minus 0.5pt}
\setlist{nosep,leftmargin=1.6em}\raggedbottom
\newtheorem{theorem}{Theorem}\newtheorem{proposition}[theorem]{Proposition}
\newtheorem{lemma}[theorem]{Lemma}
\hypersetup{pdftitle={From Dirichlet GRH to class groups and prime races},pdfauthor={OpenAI}}
\begin{document}
\thispagestyle{plain}
{\sffamily\small\color{ink} RESEARCH NOTE\hfill 25 NOVEMBER 2026}\par\vspace{7pt}
{\Huge\bfseries\color{ink} From Dirichlet GRH to class groups and prime races\par}
\vspace{3pt}
{\large An analysis of the Sernau Lean development of the Alp\"oge--Claude proof}\par
\vspace{6pt}
The repository proves a statement about every Dirichlet character. Two questions expose both its reach and its limits. How many imaginary quadratic fields have a class group of exponent six? Can one prove a persistent bias in the splitting of primes for a cubic polynomial? The first admits an explicit partial classification. For the second, the final analytic idea in the Lean proof gives an equivalence with a single degree-two Riemann hypothesis.

\heading{A partial classification with exact counts}
Let $D<0$ be a fundamental discriminant, $x=|D|$, and $P^+(x)$ the largest prime divisor of $x$. The exponent of a finite abelian group is the least positive integer that annihilates every element. Elsenhans, Kl\"uners and Nicolae (EKN) found 432 imaginary quadratic fields of exact exponent six in a census through $3.1\cdot10^{20}$ \cite{ekn}. We prove that their list is complete on the union of the two infinite families
\[
   \omega(x)\le6 \qquad\text{or}\qquad x/P^+(x)\le 10^8,
\]
where $\omega$ counts distinct prime divisors. The first family contains exactly 429 fields of exponent six. The union contains all 432 known fields. Any additional example must have at least seven ramified primes and a cofactor greater than $10^8$. The proof uses Dirichlet GRH over $\QQ$, genus theory and explicit constants.

\heading{One proved prime race and one exact open criterion}
For $p\ne23$ put $b(p)=\#\{u\bmod p:u^3-u-1=0\}-1$. Define $B(y)=\sum_{p\le y,\,p\ne23}b(p)\log p$. We prove that the Riemann hypothesis for the standard Artin $L$-function of this cubic is equivalent to a simple sign assertion
\[
                    -\int_x^{4x} B(y)\,\frac{\dd y}{y^2}>0
                    \quad\text{for every sufficiently large real }x.
\]
For the corresponding race with $b(p)=1$ on primes $1\bmod4$ and $b(p)=-1$ on primes $3\bmod4$, the repository supplies the missing zero theorem and we prove strict positivity for \emph{every} real $x>3/4$, the sharp starting point. The cubic criterion uses a bound on the total contribution of zeros derived here from the positive eta product of level 23.

\heading{Evidence and scope}
The five upstream theorem statements have independent kernel-verification records. We checked their source identity, followed the main estimate to those statements, and ran the accompanying finite computations. The deductions in this note are mathematical proofs on paper. The reported finite ranges and numerical certificates have reproducible checks. The previously published conditional theorems and computations retain their original attribution. No claim of priority is made for the deductions here.

\newpage
\section{What can be taken from the Lean repository}
The inspected revision is \href{https://github.com/lsernau/grh-lean/tree/72fead102f663d5c7a407964a13dc0cf9d977207}{\texttt{72fead102f663d5c7a407964a13dc0cf9d977207}}. Both projects pin Lean 4.35.0 and Mathlib \texttt{56a00585}. The \gh{src/Statements.lean}{public definitions} import only Mathlib. The \gh{src/Main.lean}{five theorems} give Mathlib's RH, a zeta-zero formulation, nonvanishing for all Dirichlet characters to the right of $1/2$, the critical-strip zero statement for all characters, and its primitive specialization. The nonvanishing statement allows the pole at $s=1$ for a principal character. The modulus is an arbitrary positive integer. This scope includes imprimitive characters.

Andrew Sutherland reports successful builds, axiom checks and kernel replay on both editions. The readable project has 803 local source modules, the compact one 794. The paranoid comparator accepted both. The only axioms of the five proof terms reported there are \texttt{propext}, \texttt{Classical.choice} and \texttt{Quot.sound} \cite{mit}. Our source check reproduces the published Git archive digest, the equality of the public statement files, and import reachability from \texttt{Main}. We used the independently published kernel runs. We did not run a separate Lean build or kernel replay here.

\heading{Where the square-root exponent enters}
Fix an admissible modulus $q_0$ divisible by 18480 and an auxiliary quadratic character $\chi_d$. The construction detects the centered prime-power sum
\[
 R_q(x)=\sum_{\substack{n\le x\\ n\equiv 1\pmod q}}
      \Lambda(n)(1-\chi_d(n))(1/n-1/x)-\frac{\log x}{\varphi(q)}-c_q.
\]
The \gh{src/New/Reduction/MainInequality.lean}{main inequality} compares two estimates for the same covolume, corrected by twice a lattice index. Short integral forms bound it above, a spanning estimate supplies many independent forms, and a twisted determinant bounds it below. Cancellation of the common main term gives
$m\,\mathrm{rem}\le f a(\log b)^5$ in the repository's parameters. The chain has $m\asymp\sqrt{x}$ lifts. A much larger auxiliary scale absorbs two prime-number-theorem errors. The \gh{src/New/Reduction/OneSidedBound.lean}{resulting one-sided estimate} is $R_{q_0}(x)\ll_{q_0,d,\epsilon}x^{-1/2+\epsilon}$. The \gh{src/New/Reduction/Zeros/Criterion.lean}{zero criterion} chooses a twist to isolate untwisted zeros and uses a Landau-type argument on the resulting Mellin transform. This last passage suggests the prime-race construction on pp.~\pageref{races}--\pageref{obstruction}.

\heading{Two short routes to established arithmetic targets}
\textbf{Grosswald's conjecture.} For every prime $p>409$ the least positive primitive root satisfies $g(p)<\sqrt p-2$. For $p>2791$ even the least \emph{prime} primitive root satisfies the bound. McGown, Trevi\~no and Trudgian prove both using smoothed character sums modulo $p$. Their Theorem~1 and the character expansion in their Lemma~1 argument require only $\zeta$ and Dirichlet $L$-functions \cite{grosswald}. The repository therefore discharges the hypothesis of their published proof.

\textbf{Small exponents.} For quadratic $K$ of fundamental discriminant $D$,
$\zeta_K(s)=\zeta(s)L(s,\chi_D)$, where $\chi_D$ is the primitive Kronecker character. Both factors meet the zero hypothesis in EKN's Theorem~2. The counts at exact exponents $1,2,3,4,5,8$ are respectively $9,56,17,203,27,778$ \cite{ekn}. For exponent five, the norm argument on the next page and the published small split-prime bound give $|D|\le 2.3\cdot10^{20}$, inside the published census. These are applications of previously proved implications.

\newpage
\section{The exponent-six classification}
Write $G_D=\Cl(\QQ(\sqrt D))$ and $E(D)=\exp G_D$. Set $X_0=3.1\cdot10^{20}$. EKN's census and their quadratic split-prime bound imply that any further exponent-six example has
\begin{equation}
                  X_0 < x=|D|\le 2.5\cdot10^{25}.                 \tag{1}
\end{equation}
We exploit the exponent of the subgroup of squares in $G_D$.

\begin{lemma}\label{norm}
If $p$ splits in $K=\QQ(\sqrt D)$, $\mathfrak p\mid p$, and $[\mathfrak p]^j=1$, then $x\le4p^j$. If $j=2r$ then also $P^+(x)<4p^r$.
\end{lemma}
\begin{proof}
Write $\mathfrak p^j=(\alpha)$ with $\alpha=(a+b\sqrt D)/2\in\mathcal O_K$. The valuations at $\mathfrak p$ and $\bar{\mathfrak p}$ differ, so $b\ne0$. Norms give $4p^j=a^2+xb^2$, proving the first assertion. For even $j$, the positive integers $2p^r-a$ and $2p^r+a$ have product $xb^2$. Each is strictly less than $4p^r$. Every prime factor of $x$ divides at least one of them.
\end{proof}

\begin{theorem}\label{six}
All exponent-six fields for which $\omega(x)\le6$ or $x/P^+(x)\le10^8$ occur in EKN's list. If an omitted example has exactly seven prime divisors, then $x<2.1\cdot10^{22}$. If it has exactly eight, then $x<2.2\cdot10^{24}$.
\end{theorem}
\begin{proof}
Let $t=\omega(x)$ and let $F$ be the genus field of $K$. The principal genus theorem identifies $\operatorname{Gal}(F/K)$ with $G_D/G_D^2$. The field $F$ is the compositum of the quadratic fields associated with the $t$ prime discriminants in $D$, including the appropriate prime discriminant at 2. It has degree $n=2^t$ over $\QQ$ and is unramified at all finite primes over $K$ \cite{cox}. The tower formula gives
\[
                    \log |\Delta_F|=(n/2)\log x .
\]
Every character of $F/\QQ$ gives a Dirichlet character. Hence the repository provides the zeros required by the Bach--Sorenson bound for a rational prime $p$ that splits completely in $F$ \cite{bach}. The Artin map puts $[\mathfrak p]$ in $G_D^2$ for either prime of $K$ over $p$. Since $E(D)=6$, every square in $G_D$ has cube one. Lemma~\ref{norm} now gives
\begin{equation}
                         (x/4)^{1/6}\le\sqrt p.                  \tag{2}
\end{equation}
Take $y=\log x$. In (1), $47<y<59$. For $t\le4$, Bach--Sorenson Theorem~5.1 gives $\sqrt p\le4\log|\Delta_F|+2.5n+5\le32y+45$. Their Table~3 for base field $\QQ$ gives the sharper bounds below. Its entries are triples $(a,b,c)$ for $\sqrt p\le a\log|\Delta_F|+bn+c$.

\begin{center}\small
\begin{tabular}{r r l l}\toprule
$t$ & $n$ & Table range for $\log|\Delta_F|$ & Upper bound on $\sqrt p$\\\midrule
5&32&$100$--$1000$ & $1.745(16y)+0.04(32)+5.8$\\
6&64&$1000$--$10000$ & $1.257(32y)+7.3$\\
7&128&$1000$--$10000$ & $1.257(64y)+7.3$\\
8&256&$1000$--$10000$ & $1.257(128y)+7.3$\\\bottomrule
\end{tabular}\end{center}
For any of these affine bounds $M(y)=ay+b$, $M(y)\e^{-y/6}$ decreases when $y>6$. For $t\le6$, all the bounds are at most $40.224y+7.3$ on (1). At $X_0$, $y<47.2$, so this last bound is $<1906$ and $4(1906)^6<X_0$. This contradicts (2) throughout the candidate range. The same comparison in the last two rows, at $2.1\cdot10^{22}$ and $2.2\cdot10^{24}$ respectively, proves the two asserted upper limits. Exact rational checks of the endpoints accompany the note.

\newpage
\heading{The cofactor restriction and the exact counts}
To complete the cofactor assertion of Theorem~\ref{six}, take a prime $p$ split in the quadratic field $K$ itself. EKN Theorem~3, using Bach--Sorenson, gives $\sqrt p\le M=1.881\log x+6.18$ in (1). The second part of Lemma~\ref{norm} with $j=6$ gives
\[
                    \frac{x}{P^+(x)}>\frac{x}{4M^6}.
\]
The right-hand side increases on (1). At $X_0$ one has $M<95$ and $X_0/(4\cdot95^6)>10^8$. This proves the cofactor assertion.
\end{proof}

We recomputed the fields up to $5{,}761{,}140$, the largest $|D|$ in EKN's exponent-six table. The histogram uses the number of distinct rational primes ramified in the quadratic field.
\begin{center}
\begin{tabular}{lrrrrrrr}\toprule
$\omega(|D|)$ & 2&3&4&5&6&7&Total\\\midrule
Fields of exact exponent six&60&141&140&76&12&3&432\\\bottomrule
\end{tabular}\end{center}
Thus 429 have $\omega\le6$. The other three discriminants are $-3{,}892{,}980$, $-4{,}696{,}692$ and $-5{,}761{,}140$. All 432 satisfy the cofactor restriction. The largest observed cofactor is $204{,}204$. The completeness through $X_0$ uses EKN's published computation. Our independent finite calculation covers the smaller range stated here.

\heading{Checking class groups by binary quadratic forms}
A negative fundamental discriminant gives one ideal class per reduced primitive positive definite binary form $(a,b,c)$, with $b^2-4ac=D$, $|b|\le a\le c$ and the usual nonnegative boundary sign. The programs enumerate these forms using $3a^2\le |D|$. A split prime whose small power is nonprincipal rules out a candidate efficiently. For a remaining field we test powers of every class, so large class number alone cannot produce a false positive.

Here is the power test. Change variables unimodularly until $(a,D)=1$; such a representative exists by avoiding the bad residues at primes dividing $D$. Set $\delta=D\bmod 2$, $\omega_K=(\delta+\sqrt D)/2$ and $f(T)=T^2-\delta T+(\delta^2-D)/4$. The form corresponds, up to inverse convention, to $(a,\omega_K-u)$ for $u=(\delta-b)/2\pmod a$. The derivative at $u$ is a unit mod $a$. Lift the root to $u_j\pmod {a^j}$ by Hensel's rule. The kernel of evaluation $\omega_K\mapsto u_j$ mod $a^j$ contains the $j$th ideal power and has the same norm. Reducing its form tests principality. An abelian group killed by six with order divisible by six has exponent exactly six.

For exponent five we independently enumerated the $75{,}988$ negative fundamental discriminants with $|D|\le250{,}000$. Their forms numbered $11{,}689{,}615$. Of 196 fields with nontrivial class number a power of five, exactly 27 had exponent five. Of these, 25 have class number five. At $D=-12{,}451$ and $-37{,}363$ the group is $C_5\times C_5$. EKN's analytic cutoff and far larger census justify completeness at every discriminant. The archive gives both CSV lists and the arithmetic that produced them.

\heading{What the restrictions leave to investigate}
Additional exponent-six fields would require a fairly balanced factorization of a very large discriminant with many prime factors. For seven or eight factors the last two rows on the preceding page further shorten the interval. The degree $2^t$ of the genus field explains quantitatively why the same constants lose strength as the number of factors increases.

\newpage
\section{A prime race in a fixed window}\label{races}
Write $b(p)$ for either of the following two weights. In the quadratic case omit $p=2$. In the cubic case omit $p=23$.
\begin{center}
\begin{tabular}{lll}\toprule
Case & Prime weight $b(p)$ & Corresponding $L$-function\\\midrule
Modulo 4 & $\chi_4(p)$, taking values $1,-1$ & Dirichlet beta $\beta(s)$\\
Cubic of discriminant $-23$ & $\#\{u\in\mathbb F_p:u^3-u-1=0\}-1$ & standard Artin $L(s,\rho)$\\\bottomrule
\end{tabular}\end{center}
Define $B(y)=\sum_{p\le y} b(p)\log p$. For a fixed dilation $h>1$, let
\begin{equation}
 J_h(x)=-\int_x^{hx} B(y)\,\frac{\dd y}{y^2}
       =-\sum_{p\le hx}b(p)\log p\left(\frac{1}{\max(p,x)}-\frac1{hx}\right).       \tag{3}
\end{equation}
Thus every value of $J_h$ involves finitely many primes. A positive value in the modulo-four race favors primes $3\bmod4$. In the cubic race it favors irreducible primes over twice the completely split primes. The one-root primes carry weight zero.

\begin{theorem}\label{race}
For the modulo-four race, $J_4(x)>0$ for every real $x>3/4$, $J_4(x)=0$ for $0<x\le3/4$, and $\liminf_{x\to\infty}\sqrt x\,J_4(x)>3/8$.
For the cubic, the following are equivalent:
\begin{enumerate}[label=\textup{(\roman*)}]
\item Every zero of $L(s,\rho)$ in the critical strip has real part $1/2$.
\item $J_4(x)\ge0$ for every sufficiently large real $x$.
\item $\displaystyle\liminf_{x\to\infty}\sqrt x\,J_4(x)\ge1-\frac{9}{\pi\sqrt{23}}>0.4026$.
\end{enumerate}
The same equivalence holds with strict positivity in \textup{(ii)}.
\end{theorem}

\heading{Why the square terms have the decisive sign}
For the cubic, the polynomial is irreducible by the rational root test and has nonsquare discriminant $-23$. Its splitting field has group $S_3$. The standard representation $\rho$ has character equal to the number of fixed letters minus one. For unramified $p$ let $g_p$ denote Frobenius. The traces needed for the prime-power sum are as follows.
\begin{center}\small
\begin{tabular}{lrrrr}\toprule
Cycle type & Number in $S_3$ & Roots & $\mathrm{tr}\,\rho(g_p)$ & $\mathrm{tr}\,\rho(g_p^2)$\\\midrule
Identity &1&3&2&2\\
Transposition &3&1&0&2\\
Three-cycle &2&0&$-1$&$-1$\\\bottomrule
\end{tabular}\end{center}
Let $A(y)=\sum_{p^k\le y}\mathrm{tr}\,\rho(g_p^k)\log p$, using $\chi_4(p)^k$ in the quadratic case. The mean of the last column is $(2+6-2)/6=1$. Chebotarev's theorem and partial summation give a square contribution $\sqrt y+o(\sqrt y)$ for the fixed cubic field, without any hypothesis on its zeros. For $\chi_4$ the squares have weight one at every odd prime, so the ordinary prime number theorem suffices. Powers with $k\ge3$ contribute $O(y^{1/3}\log^2y)$ in both cases. With $H_h=-\int_x^{hx}A(y)y^{-2}\dd y$ this gives
\begin{equation}
 \sqrt x\,[J_h(x)-H_h(x)]\longrightarrow 2(1-h^{-1/2}).                     \tag{4}
\end{equation}
Passing from primes to prime powers changes the main term at precisely the scale of the desired bias.

\newpage
\heading{Bounding the oscillation of the prime powers}
Use $L_*$ for the $L$-function with the omitted local factor removed, and $D_*(s)=-L_*'(s)/L_*(s)$. In the cubic case $L_*=(1-23^{-s})L$ \cite{serre}. Write $a_{p^k}=\mathrm{tr}\,\rho(g_p^k)\log p$, or $a_{p^k}=\chi_4(p)^k\log p$. The factor $1/n-1/x$ is the same smoothing that appears in the repository. Set
\[
 S(x)=\sum_{p^k\le x}a_{p^k}(p^{-k}-x^{-1}),\qquad H_h(x)=S(x)-S(hx).
\]
Suppose now that RH holds for this individual $L$. Denote its nontrivial zeros, repeated with multiplicity, by $\rho'=1/2+i\gamma$, and put $Z=\sum_{\rho'}(1/4+\gamma^2)^{-1}$. The smoothed explicit formula is
\begin{equation}
 S(x)=D_*(1)+\sum_{\rho'}\frac{x^{\rho'-1}}{\rho'(1-\rho')}+O(x^{-3/4}).       \tag{5}
\end{equation}
Indeed the inverse Mellin kernel for the sum is $D_*(s)x^{s-1}/[s(s-1)]$ on $\Re s=2$. Move the line to $1/4$. The pole at 1 comes from the kernel, and the residue of $D_*$ at a zero of multiplicity $m$ is $-m$. The omitted local factor is nonzero and has bounded logarithmic derivative on this new line. The remaining vertical integral is $O(x^{-3/4})$: the standard partial fraction formula for the logarithmic derivative of a fixed completed $L$-function, Stirling's formula and the $O(\log(2+T))$ count of zeros in a unit interval give $D_*(1/4+it)=O(\log^2(3+|t|))$. For the horizontal segments take heights separated by $\gg1/\log T$ from nearby zero ordinates. The same bound there and the squared denominator let the segments tend to zero. These fixed-function estimates and their Artin version are described in \cite{ng}. The series in (5) converges absolutely, since the number of ordinates up to $T$ is $O(T\log T)$.

The constant $D_*(1)$ disappears from $H_h$. Equations (4) and (5) yield
\begin{equation}
 \liminf_{x\to\infty}\sqrt x\,J_h(x)
       \ \ge\ 2(1-h^{-1/2})-(1+h^{-1/2})Z .                        \tag{6}
\end{equation}
For the completed primitive $L$-function $\Li$, both examples have sign $+1$ and $\Li(1/2)\ne0$. The Hadamard product for the even entire function $z\mapsto\Li(1/2+z)$ pairs zeros at $\pm i\gamma$. Its order is one, and evenness removes the linear exponential. Differentiating the paired product at $z=1/2$ gives
\begin{equation}
                          Z=2\,\Li'(1)/\Li(1).                      \tag{7}
\end{equation}

\heading{The modulo-four race}
Here $\Li(s)=(4/\pi)^{(s+1)/2}\Gamma((s+1)/2)\beta(s)$, and the positive alternating beta series gives $\beta(s)>0$ for all real $s>0$. Formula (7) gives
\[
 Z=\log(4/\pi)-\gamma_{\!E}+2\beta'(1)/\beta(1)<5/12 .
\]
Indeed $\beta(1)=\pi/4>3/4$. The alternating upper bound
$\beta'(1)\le\sum_{j=1}^{15}(-1)^{j+1}\log(2j+1)/(2j+1)<1/4$
uses terms decreasing from 3 onwards. The finite inequality is verified by rational logarithm enclosures in the archive. Use also $\log(4/\pi)<1/4$ and $\gamma_E>1/2$. Dirichlet GRH supplies RH for $\beta$. Taking $h=4$ in (6) proves the asymptotic lower bound $1-\tfrac32 Z>3/8$.

\newpage
\heading{Positivity at every scale for the modulo-four race}
To reach the sharp endpoint $x=3/4$, we need an explicit large-$x$ estimate and a finite proof covering the remaining real interval. For $x>1$ shift the contour defining $S(x)$ farther to the left for $L=\beta$. In addition to the constant $D_*(1)$ and the zero series from (5), the exact formula has terms
\[
 \frac{\beta'(0)}{x\beta(0)}
 -\sum_{j=0}^{\infty}\frac{x^{-2j-2}}{(2j+1)(2j+2)}.
\]
These are the residues at $s=0$ and the negative odd integers, respectively. To justify the full shift, move to $\Re s=-2M$ along good horizontal heights and then let $M\to\infty$. On that line the functional equation, the digamma reflection formula, Stirling's bound and the convergent logarithmic derivative at $1-s$ give $O(\log(2+M+|\Im s|))$ for the logarithmic derivative, uniformly for integers $M\ge1$. The resulting integral is $O(x^{-2M-1}\log(2+M)/M)$. The functional equation also gives
\[
 \beta'(0)/\beta(0)=\gamma_E+\log(\pi/2)-\beta'(1)/\beta(1)>0.
\]
The trivial-zero coefficients sum to $\log 2<1$. Consequently
$H_4(x)=S(x)-S(4x)\ge-3Z/(2\sqrt{x})-x^{-2}$.

The higher prime powers supply an explicit gain. For $p\equiv3\bmod4$, exponents $k\ge2$ contribute alternately $+\log p,-\log p,\ldots$, so their sum is nonnegative; for $p\equiv1\bmod4$ each term is positive. Every odd prime in $t^{1/3}<p\le\sqrt t$ contributes $\log p$. Therefore for $t\ge65{,}536$,
\[
 A(t)-B(t)\ \ge\ \vartheta(\sqrt t)-\vartheta(t^{1/3})
 \ >\ 0.84\sqrt t-1.02\,t^{1/3}\ >\ 0.67\sqrt t .
\]
Here $\vartheta(u)=\sum_{p\le u}\log p$. Rosser--Schoenfeld, Theorems 9--10, give the upper bound $1.01624u$ for all positive $u$ and the lower bound $0.84u$ for $u\ge101$ \cite{rs}. We used $t^{1/6}>6$. Integration yields
\[
 \sqrt x\,J_4(x)>0.67-5/8-x^{-3/2}>0 \qquad(x\ge65{,}536).
\]

On $[1,65{,}536]$, (3) is continuous and affine in $1/x$ between successive breakpoints $p$ and $p/4$ for odd primes $p$. It suffices to check these breakpoints and the two endpoints. The included \texttt{mod4\_all\_x.py} encloses each $\log p$ using integer arithmetic. Reduce the argument to $1\le a/b\le2$, put $w=(a-b)/(a+b)\le1/3$, and round outward in
\[
 \log(a/b)=2\sum_{j=0}^{15}\frac{w^{2j+1}}{2j+1}+R,\qquad
 0\le R\le\frac{2(1/3)^{33}}{33(1-1/9)} .
\]
At scale $2^{48}$ the tail costs less than one integer unit. The certificate covers all $29{,}541$ test points and gives $J_4(x)\ge3681/10^6$ on the whole interval. Finally, for $3/4<x<1$ the only contributing prime is 3, with weight $\log3(1/3-1/(4x))>0$. For $0<x\le3/4$ no term has positive support. This proves the exact starting point of the race.
 
\newpage
\section{Controlling a degree-two zero sum}
For $K=\QQ[\alpha]$, $\alpha^3-\alpha-1=0$, Serre's description identifies the standard Artin factor with the Hecke $L$-function of the weight-one form
\[
 f(z)=\eta(z)\eta(23z),\qquad
 \zeta_K(s)=\zeta(s)L(s,\rho).                                      \tag{8}
\]
The completed function is $\Li(s)=(\sqrt{23}/(2\pi))^s\Gamma(s)L(s,\rho)$, with $\Li(s)=\Li(1-s)$ \cite{serre,zagier}. The eta product is positive on the positive imaginary axis. Its Mellin integral makes $\Li(s)>0$ for all real $s$, including $s=1/2$. We need quantitative control of its logarithmic derivative as well.

\begin{proposition}\label{eta}
Let $F(y)=\prod_j\eta(i d_j y)^{r_j}$ with finitely many $d_j,r_j>0$. For the entire Mellin transform $M(s)=\int_0^\infty F(y)y^s\dd y/y$ let
\[
 c=\frac{\pi}{12}\inf_{y>0}\sum_j r_j\max\{d_j y,(d_j y)^{-1}\}.
\]
For every real $\sigma$, $(\log M)''(\sigma)\le1/c$. The transform has no zero at any $s$ with $|\Im s|\le(\pi/2)\sqrt c$.
\end{proposition}
\begin{proof}
Set $v=\log y$. The product formula for $\eta(i\e^v)$ starts with $\exp(-\pi\e^v/12)$. For $z=a\e^v$, the derivative of $\log(1-\e^{-z})$ is $z/(\e^z-1)$, decreasing in $v$. Thus $-(\log\eta(i\e^v))''\ge(\pi/12)\e^v$. Termwise differentiation is justified locally by exponential convergence. Eta inversion gives $\eta(i\e^{-v})=\e^{v/2}\eta(i\e^v)$, making this second logarithmic derivative even. Apply the two lower bounds to each factor separately. The potential $V(v)=-\log F(\e^v)$ satisfies $V''\ge c$. The eta product and its inversion also give double exponential decay in $v$ at both ends, up to an exponential linear factor. The integral is entire and all differentiations below are justified.

Give $U$ density proportional to $\exp(\sigma v-V(v))$, and write $\mu=\mathbb E U$. Logarithmic differentiation gives $(\log M)''(\sigma)=\operatorname{Var}U$. Integration by parts gives $\mathbb E[(U-\mu)(V'(U)-\sigma)]=1$. Subtracting $V'(\mu)$ in place of $\sigma$ leaves this expectation unchanged. Strong convexity then gives $1\ge c\operatorname{Var}U$. Finally, $\cos w\ge\pi/4-w^2/\pi$ on $\RR$, strictly except at $w=\pm\pi/2$. On $[0,\pi/2]$ this follows by differentiating and using $\sin w\ge2w/\pi$; past $\pi/2$, $2w/\pi\ge1$. The centered characteristic function consequently has strictly positive real part whenever $t^2\operatorname{Var}U\le\pi^2/4$. It equals $\e^{-it\mu}M(\sigma+it)/M(\sigma)$.
\end{proof}

For $d_1=1,d_2=23$ the infimum is $2\sqrt{23}$, attained at $y=1/\sqrt{23}$, by comparison with $1/y+23y$. Thus $c=\pi\sqrt{23}/6$. Multiplying $M(s)$ by $23^{s/2}$ gives $\Li(s)$ and leaves the second logarithmic derivative unchanged. The symmetry forces $(\log\Li)'(1/2)=0$. Integrating the variance bound up to $1$ gives, under the cubic RH, $Z\le1/c$. Equations (6)--(7) give precisely $1-9/(\pi\sqrt{23})$ in Theorem~\ref{race}. Proposition~\ref{eta} also yields, without the cubic RH, a zero-free band of height $2.48915\ldots$ for the completed function. Earlier computations reach much farther for this particular field \cite{ng}. The proposition gives one formula for all positive eta products; no optimality or priority is asserted.

\newpage
\section{Why a sign would prove the cubic RH}\label{obstruction}
Extend (3) to all $x>0$; it vanishes for $x$ sufficiently small. For $\Re s>1$, substitution $y=xv$ and absolute convergence give
\begin{equation}
 \int_0^\infty J_h(x)x^{-s}\dd x
   =\frac{1-h^{s-1}}{s(s-1)}\,
      \sum_{p\ne23}\frac{b(p)\log p}{p^s}.                         \tag{9}
\end{equation}
The sum over primes equals $D_*(s)$ minus the contribution of $p^k$, $k\ge2$. That contribution is holomorphic on $\Re s>1/2$ by absolute convergence. Each zero of $L$ with $1/2<\Re s<1$ therefore gives a pole of (9), because $1-h^{s-1}$ cannot vanish there. At $s=1$ the apparent singularity of the prefactor is removable. Positivity of the eta Mellin integral gives no real zeros of $L$ on $(1/2,\infty)$, so (9) is regular on that real interval.

Suppose $J_h$ is eventually nonnegative. Landau's theorem says that the Mellin transform of a nonnegative tail has a singularity at any finite real abscissa of convergence. For completeness, expand its transform at a real point just to the right of the abscissa, in the direction of decreasing $s$. Its Taylor coefficients are the integrals of the tail against nonnegative powers of $\log x$ divided by factorials (start the tail at $x\ge1$). Analyticity at the abscissa would permit a small disk of continuation there. Expanding at a nearby point to its right gives a Taylor radius extending to the left of the abscissa. Monotone convergence of the exponential series would then give convergence of the original integral at that point, a contradiction. Removing a finite interval contributes an entire function in (9).

Since there is no real singularity above $1/2$, the integral is holomorphic in the whole right half-plane $\Re s>1/2$. All the poles just identified are excluded. The functional equation reflects any nontrivial zero to the left of $1/2$ into the right half. This proves (ii)$\Rightarrow$(i); (i)$\Rightarrow$(iii) was proved on the previous pages, and (iii)$\Rightarrow$(ii) is immediate. An eventually nonpositive race would also yield RH by applying Landau to $-J_h$; for $h=4$ that would contradict the positive lower bound. In particular, the cubic race takes positive values arbitrarily far out even without RH.

\heading{The precise obstruction to a congruence substitution}
Let $E$ be the splitting field of the cubic and $m$ any fixed modulus. For every reduced residue class $r\bmod m$, the prime average of $b(p)$ conditional on $p\equiv r\bmod m$ is zero. Indeed $E\cap\QQ(\zeta_m)$ is abelian over $\QQ$, so it is either $\QQ$ or the unique quadratic subfield of $E$. Chebotarev in the compositum makes Frobenius equidistributed, with element multiplicities, on $S_3$, on $A_3$ or on its odd coset as appropriate. The trace $2,0,-1$ has mean zero on each. Its unconditional mean square is $(4+0+2)/6=1$.

Thus for every real function $w$ on the units mod $m$,
\[
 \lim_{X\to\infty}\frac1{\pi(X)}
 \sum_{\substack{p\le X\\p\nmid 23m}}(b(p)-w(p\bmod m))^2
 =1+\frac1{\varphi(m)}\sum_{r\in(\mathbb Z/m)^\times}w(r)^2 .
\]
The best fixed periodic approximation is zero in this precise norm. A cubic divisibility or height estimate would have to capture information that no fixed collection of Dirichlet-character weights captures. The statement allows other constructions, including moduli varying with the cutoff. Neither the prime-race reformulation nor the numerical examples supply this new estimate.

\newpage
\section{Checks, provenance and further questions}
For the cubic, $p=2,5,59$ respectively exhibit the irreducible, mixed and fully split cases. The program counts roots by $\deg\gcd(X^3-X-1,X^{p^k}-X)$ and checks traces against Frobenius powers. It independently verifies the difference identity $H_h(x)=S(x)-S(hx)$. The two final columns are rounded values of (3). Their computation illustrates the scale and sign; the proofs do not extrapolate the table.
\begin{center}\small
\begin{tabular}{rrrrrr}\toprule
$x$ & Three roots&One root&No roots&$\sqrt{x}\,J_4$ cubic &$\sqrt{x}\,J_4$ mod 4\\\midrule
100&1&14&9&1.343855&0.674377\\
1000&25&87&55&0.9648&0.747457\\
10000&193&627&408&1.150540&0.961868\\
100000&1567&4823&3201&1.1356&1.032079\\\bottomrule
\end{tabular}\end{center}

\heading{Further research questions}
\begin{enumerate}
\item Can the genus-field argument be sharpened for seven or more ramified primes by splitting in a smaller subfield and multiplying a prime ideal by ramified ideals of small norm? The competing quantities are the reduction in degree and the increase in the generator norm. Their effect depends on the actual rank of the ramified classes in the genus quotient.
\item Can a theta-lift construction produce an upper bound for a smoothing involving the difference of the two form classes of discriminant $-23$? The principal form $u^2+uv+6v^2$ detects completely split primes and $2u^2+uv+3v^2$ detects irreducible ones \cite{serre}. Equation (9) specifies the analytic output to aim for.
\item Which other self-dual Artin or modular $L$-functions have a prime-square bias large enough to dominate their zero sum? The inequality in (6) and the curvature calculation in Proposition~\ref{eta} give verifiable sufficient conditions. Suzuki's sign criteria and Ng's computations give relevant precedents \cite{suzuki,ng}.
\end{enumerate}

\heading{Reproducibility}
The file \texttt{README.md} gives commands. The \texttt{checks/} directory contains the pinned source audit, two finite class-group enumerations and CSV files, rational endpoint certificates, and prime-race examples. Finite class tests are exact; the sample weighted prime sums use floating point. The global inferences rely on the upstream theorem, the cited published analytic results and the published census. The new arguments here have not undergone Lean kernel checking.

\newpage
\begin{thebibliography}{12}\normalsize\setlength{\itemsep}{0pt}\setlength{\parskip}{0pt}
\bibitem{repo} L. Sernau and collaborators, \href{https://github.com/lsernau/grh-lean}{\emph{GRH for Dirichlet $L$-functions, in Lean 4}} (2026), based on \href{https://alpo.ge/remark.pdf}{L. Alp\"oge's announcement}. See the \gh{GUIDE.md}{reading guide}.
\bibitem{mit} A. Sutherland, \href{https://math.mit.edu/~drew/shortgrh/}{\emph{Verification files for the shorter Lean proof of RH and GRH}}. See also \href{https://mrubinst.github.io/grh-lean-compact-check/}{M. Rubinstein's compact-edition check}.
\bibitem{ekn} A.-S. Elsenhans, J. Kl\"uners and F. Nicolae, \href{https://arxiv.org/abs/1803.02056}{\emph{Imaginary quadratic number fields with class groups of small exponent}}, Acta Arith. 193 (2020), 217--233, Theorems 1--3, Tables 1--2. \href{https://math.uni-paderborn.de/en/ag/ca/research/exponent/exponent0}{Authors' tables}.
\bibitem{bach} E. Bach and J. Sorenson, \href{https://doi.org/10.1090/S0025-5718-96-00763-6}{\emph{Explicit bounds for primes in residue classes}}, Math. Comp. 65 (1996), 1717--1735, Theorem 5.1 and Table 3.
\bibitem{cox} D. A. Cox, \emph{Primes of the Form $x^2+ny^2$}, 2nd ed., Wiley, 2013, \S6 (genus theory for field discriminants).
\bibitem{grosswald} K. McGown, E. Trevi\~no and T. Trudgian, \href{https://arxiv.org/abs/1508.05182}{\emph{Resolving Grosswald's conjecture on GRH}}, Funct. Approx. Comment. Math. 55 (2016), 215--225.
\bibitem{serre} J.-P. Serre, \href{https://www.math.mcgill.ca/goren/SeminarOnCohomology/Serre.pdf}{\emph{On a theorem of Jordan}}, Bull. Amer. Math. Soc. 40 (2003), 429--440, \S5.3 and notes.
\bibitem{zagier} D. Zagier, \href{https://www.its.caltech.edu/~matilde/Zagier123ModularForms.pdf}{\emph{Elliptic Modular Forms and Their Applications}}, in \emph{The 1-2-3 of Modular Forms}, 2008, Prop. 14.
\bibitem{ng} N. Ng, \href{https://www.cs.uleth.ca/~nathanng/RESEARCH/phd.thesis.pdf}{\emph{Limiting Distributions and Zeros of Artin $L$-Functions}}, PhD thesis, 2000, Chs. 3--5. The conductor-23 computations appear in \S4.3.5.
\bibitem{suzuki} M. Suzuki, \href{https://arxiv.org/abs/2411.07436}{\emph{On variants of Chebyshev's conjecture}}, Ramanujan J. 68 (2025), article 95.
\bibitem{rs} J. B. Rosser and L. Schoenfeld, \href{https://doi.org/10.1215/ijm/1255631807}{\emph{Approximate formulas for some functions of prime numbers}}, Illinois J. Math. 6 (1962), 64--94, Theorems 9--10.
\end{thebibliography}
\end{document}
